\section{Discusión}
\label{sec:discusion}

\subsection{El hallazgo tiene dos capas}

El resultado principal es que un umbral de \texttt{confidence} no transfiere entre implementaciones y que el
promedio no describe el fenómeno. Al mirar celda a celda aparecen dos capas distintas, y confundirlas lleva a conclusiones opuestas (documentado en \texttt{docs/erratas\_especificacion.md}, E2).

\textbf{Capa 1, la escala.} Dos definiciones pueden ordenar las decisiones igual y aun así estar en escalas distintas. $C_{\max}=p$ y $C_{\mathrm{Jev}}=(n\cdot p-1)/(n-1)$ inducen el mismo umbral en orden, de modo que existe un mapa cerrado $\psi=\phi^{-1}$ entre ellas. El número no viaja ($|\Delta|$ mediano de 0.076 en un sentido del par \texttt{pmax} $\leftrightarrow$ \texttt{jev} y 0.257 en el otro, nivel \texttt{raw}), pero el mapa sí transfiere (0.0013 de mediana, 0.0049 de media): la escala se corrige, la cifra no sirve tal cual.

\textbf{Capa 2, la forma.} Entre definiciones de clases distintas ($p_{\max}$, entropía, margen) no existe ni número ni mapa. \texttt{pmax} $\to$ \texttt{ent} tiene mediana de 0.77: es el caso en que la confianza «alta» de una definición y la de la otra ordenan las decisiones de forma diferente, no solo con otra escala. Aquí no hay nada que traducir y el porte directo es lisa y llanamente una extrapolación inválida.

La coexistencia de las dos capas es la razón de que el resultado no se pueda resumir con una media: la misma tabla contiene identidades exactas (14\,\% de las celdas con $n=2$ transfieren bien) y desastres (celdas con $\Delta=0.80$ y cobertura 0).

\subsection{Por qué el orden se preserva y el número no}

Sea $f$ una transformación creciente de $p_{\max}$. Un umbral $\tau$ sobre $f(p_{\max})$ selecciona el mismo conjunto de decisiones que el umbral $f^{-1}(\tau)$ sobre $p_{\max}$, porque $f$ monótona no cambia el orden. Por eso $C_{\mathrm{Jev}}$ y $C_{\max}$ (fijado $n$) seleccionan lo mismo con el umbral traducido; $C_{\mathrm{mar}}$ también, pero \textbf{solo a $n = 2$}, donde vale $2\cdot p_{\max}-1$: a $n \ge 3$ el margen depende de las dos probabilidades mayores y no es función de $p_{\max}$. Pero el valor de $f$ depende de $n$ y de la definición, así que la misma cifra $0.7$ significa coberturas distintas. Y cuando $f$ deja de ser una función de $p_{\max}$ (la entropía usa toda la distribución, no solo el máximo), ni siquiera existe el mapa.

Esto es exactamente lo que predice la Proposición 1 y lo que mide el factor 109$\times$: el costo de confundir «mismo orden» con «mismo número».

\subsection{El umbral es una política, no un número suelto}

En un sistema que decide, el umbral no es un parámetro cosmético: define cuándo el modelo actúa y cuándo se abstiene. Un umbral trasladado sin mapa no solo cambia una métrica de calibración, cambia la política de abstención. Las celdas de \S5.6 lo muestran en su forma extrema: con $q=0.80$ puede terminar actuando en el 0\,\% de los casos. Para quien integra un modelo tipo Jev, la conclusión práctica es que el campo \texttt{confidence} no es un contrato entre implementaciones: es un número local, válido dentro de la definición que lo produjo.

Y sin embargo el ecosistema lo trata como si lo fuera. La revisión de implementaciones (Kev, Laya, AnyJev, open-decision-ai) muestra que cada una publica su fórmula bajo el mismo nombre y que ninguna documenta la relación con las demás. Ese vacío es el que este instrumento llena: no descubre que hay sobre-confianza (eso ya se midió puertas adentro de varias implementaciones), sino que cuantifica el error de trasladar el umbral entre ellas y entrega el mapa cuando existe.

\subsection{Calibración de las implementaciones medidas}

La literatura reciente mide la calibración dentro de una definición (por ejemplo, ECE de Laya tal como se distribuye). Este trabajo mide entre definiciones, que es una pregunta distinta y hasta ahora sin tabla. Los dos resultados no se contradicen: un modelo puede estar bien calibrado respecto de su propia definición y aun así ser inutilizable como fuente de umbral para otra. La cifra de \S5.7 (el ECE cambia 2.7$\times$ según la definición, sobre el mismo dato) ilustra que buena parte de las discrepancias que se atribuyen al modelo son, en realidad, discrepancias del medidor.

\subsection{Limitaciones}

\begin{itemize}
\item Cobertura de modelos. Cinco configuraciones de hasta 0.8B, todas las que caben en 6 GB. Los modelos de 4B a 9B del ecosistema quedan como extensión declarada.
\item Dos grupos ausentes. \texttt{clinc150} en el checkpoint base de Laya (ventana de 512 tokens, 126 de 151 opciones) y \texttt{kev-suites} (no carga en el Hub). Ninguno altera las conclusiones, pero acotan el alcance medido a 23 de 25 combinaciones.
\item El mapa existe dentro de clases y a cardinalidad fija. Las celdas sin mapa (2\,520 de 3\,564) no son «fallos»: son el resultado de que no haya traducción matemática posible. La matriz las marca como tales.
\item El control crudo infla el porte directo en las celdas con empates de umbral (modelo sin entrenar, salidas saturadas). Por eso la comparación limpia excluye las celdas dominadas por empates y la de referencia es la de 312 celdas (109$\times$, nivel \texttt{raw}).
\item Las definiciones opacas (\texttt{act\_probability}, \texttt{open-decision-ai}) quedaron fuera de la matriz por decisión de alcance, ya que no son funciones de la distribución.
\end{itemize}

\subsection{Relación con los trabajos previos}

Este resultado es un instrumento, y se sostiene aunque otro publique un hallazgo cercano. La calibración de un modelo puntual (por ejemplo, la medición de Laya tal cual se distribuye) es un hallazgo fechado; la matriz de transferencia y la librería que la reproduce son infraestructura. El instrumento no compite con esa literatura: la cita, la usa como punto de partida y la extiende al eje que le faltaba, el eje entre definiciones.

Las demás aristas del eje (por qué el término de RL no induce calibración, la ablación del objetivo de entrenamiento, la cartografía del ecosistema, la coherencia frente a la negación) están ocupadas por trabajos recientes y quedan fuera de esta tesis. Aquí lo medido es, deliberadamente, el hueco que quedaba abierto: la transferencia de umbrales entre definiciones incompatibles del mismo campo.
